\section{Introduction}

One of the most robust findings in labor economics is the \textit{child penalty}: the sharp drop in female earnings that follows childbirth, with no corresponding drop for fathers. Using Danish register data, \citet{kleven2019children} document that mothers' earnings fall by roughly 20 percentage points relative to childless women in the years immediately following their first birth, while fathers are essentially unaffected. A central mechanism behind this asymmetry is the division of parental leave: in Denmark, mothers take approximately 80 percent of the total parental leave period despite a universal entitlement available to both partners.

This pattern raises a natural policy question: could mandatory paternity leave quotas --- which reserve a non-transferable block of leave specifically for fathers --- reduce the gender wage gap by redistributing the career cost of children? Several Nordic countries have introduced such quotas. Norway's ``daddy quota,'' introduced in 1993, is widely credited with significantly increasing fathers' leave-taking \citep{patnaik2019reserving}. Yet the long-run effects on wages, labor supply, and household welfare depend on mechanisms that are difficult to isolate with reduced-form methods alone: human capital depreciation during leave, behavioral responses in hours worked, and the general equilibrium interaction between both partners' careers.

\subsection{Research Question}

This paper asks: \textit{How do mandatory paternity leave quotas affect the life-cycle gender wage gap, female labor supply, and household welfare?}

Specifically, we aim to evaluate three policy designs:
\begin{enumerate}
    \item \textbf{Baseline} ($\bar{q} = 0$): no paternity mandate; the household freely chooses who takes leave.
    \item \textbf{Partial quota} ($\bar{q} = 0.1$): fathers must take at least 10 percent of total leave.
    \item \textbf{Equal sharing} ($\bar{q} = 0.5$): mandatory 50/50 split of parental leave.
\end{enumerate}

\subsection{Motivation}

\paragraph{Empirical motivation.}
The child penalty literature establishes three stylized facts that motivate this proposal. First, the earnings gap between mothers and childless women grows monotonically after first birth and remains large 10 years later \citep{kleven2019children}. Second, fathers' earnings trajectories are unaffected by childbirth. Third, cross-country comparisons reveal that countries with higher paternity leave take-up exhibit smaller child penalties \citep{kleven2019childpenalties}. These patterns suggest that the \textit{asymmetric} allocation of leave --- not childbirth per se --- is the proximate driver of the gap.

\paragraph{Theoretical motivation.}
Reduced-form estimates of leave reforms are informative but face two challenges. First, they capture only the intention-to-treat effect of a policy change for a particular cohort; a structural model can simulate any counterfactual quota level. Second, leave policy operates through multiple channels simultaneously: (i) the direct human capital cost of being out of the labor market, (ii) behavioral responses in market hours by both partners, and (iii) income effects from the replacement rate of leave benefits. A life-cycle model that integrates all three channels is needed to compute welfare effects and trace the full career trajectory of the gender wage gap.

Our model builds directly on the dynamic labor-fertility framework from lecture 5 \citep{kleven2019children}, extended to a two-earner household as in \citet{borella2023marriage}. The key innovation is a parental leave decision layer: when a child arrives, the couple chooses how to allocate leave between them, subject to a policy constraint.
